\section{Introduction}


\subsection{Background: VED and IFGT as theoretical infrastructure}


本論文は単独で完結する構造論であるが、その記述語彙と動力学的構文は、先行する二つの枠組みに依拠する。

第一の枠組みは \textbf{Vortical Enclosure Dynamics (VED)} である \citep{VEDVol1,VEDVol2,VEDVol3,VEDVol4}。VED は単一の公理「差がある」から出発し、時間・空間・標準模型・重力・宇宙論を、世界が「閉じようとして閉じない」過程として記述する理論系である。本論文においては、VED が提供する \textbf{差分発展構造} ── すなわち何らかの非閉包状態がそれ自身の動力学によって変形を続ける構造 ── を、知性現象の記述に投影する。

第二の枠組みは \textbf{Information Field Geometry Theory (IFGT)} である \citep{IFGT}。IFGT は閉包度 $C$、情報密度 $I$、情報流 $J$、情報場ポテンシャル $\Phi$ という最小語彙の組によって、閉包未完の運動を記述する語彙系である。本論文は IFGT を \textbf{完成された説明理論} としてではなく、規範性を密輸入しない \textbf{記述語彙} として用いる(詳細は \S{}4)。ここでの $I$、$J$、$\Phi$ は一般 IFGT 変数そのものではなく、知性層の観測窓に投影された有効変数である。

これら二系の選択には積極的理由がある。知性をめぐる既存語彙 ── 「能力」「目的」「意味」「判断」「正しさ」── は、いずれも何らかの規範性または主体性を密輸入する。VED と IFGT は、こうした密輸入を最小化するために設計された語彙系であり、本論文の構造論的議論にとって、現時点で最も摩擦の少ない記述基盤を提供する。

なお、本論文の中心主張(\S{}3 の三公理)は VED と IFGT に \textbf{論理的には依存しない}。三公理は他の任意の動力学的語彙でも書くことが可能である。VED と IFGT は、これら公理を \textbf{規範性なしに記述するための媒体} として採用される。

本シリーズには本論文と並行する別系統の論文が存在する。\textbf{Bio-IFGT} \citep{BioIFGT} は生命層における準閉包構造(quasi-closure regime)を扱い、本論文の四層対応(\S{}9.4)における生命層の記述語彙を提供する。\textbf{意識編} \citep{Consciousness} は本論文と同じ VED×IFGT 基盤のもとで、個体内における fast/slow layer 間の log referencing 構造を扱う。本論文の知性層の local horizon-form(\S{}9)と意識編の consciousness window は、同一の差分発展構造の異なる射影として位置づけられる。ただし、ここでいう horizon は複数の絶対的境界を意味しない。それは \DH{} が各観測窓に局所的 horizon-form として現れることを指す。これら並行論文との関係は \S{}9.4 および \S{}10.3 で詳しく扱う。

\subsection{The intelligence question: why neither ability nor knowledge suffices}


知性とは何かという問いは、伝統的に二つの大きな路線によって答えられてきた。一つは知性を \textbf{能力} とする路線、もう一つは知性を \textbf{知識} とする路線である。両者はいずれも構造的に行き詰まる。

\textbf{能力説の困難}は、能力の対象規定にある。「何かを行う能力」と言うとき、その「何か」を規定する瞬間に循環が生じる ── 知性的に行うべきことを規定する基準そのものが知性に依存する。問題解決能力、適応能力、抽象化能力など、いかなる候補を挙げても、その候補を選別する基準自体が知性の前提となる。

\textbf{知識説の困難}は、Gettier (1963) 以来よく知られている \citep{Gettier1963}。「真であり、信じられており、正当化されている命題の所有」を知識とする古典的定義は、無限の反例を生む。修正条件をいくら加えても新たな反例が生成され、これは Gettier 以来六十年以上にわたって構造的に持続している。

両者の困難は別個のものではない。両者に共通する構造的欠陥は、知性を何らかの \textbf{停止構造} ── すなわち静的に保持・参照・所有できる対象 ── として捉えようとしている点にある。能力は「持っている/持っていない」の停止判定を要請し、知識は「成立している/成立していない」の停止判定を要請する。

本論文の問題提起は、この前提そのものを反転させることである:

\begin{quote}
知性は停止構造ではなく、\textbf{停止構造を生み出し続ける運動} である。
\end{quote}

この反転は単なる視点の変更ではない。停止構造としての知性論が必然的に行き詰まる以上、運動としての知性論への移行は構造的要請である。

この立場は、本論文の枠組みからは独立に、生物物理学の側から既に到達されている。Nakagaki et al. (2000) 以来の \emph{Physarum polycephalum} 研究は、粘菌に知能を帰属させずに、純粋に動力学として知性的振る舞いを記述する立場を一貫して採ってきた \citep{Nakagaki2000,Tero2007}。中垣自身の表現「粘菌は知能を持たないが、知能的に振る舞う」は、本論文の中心立場 ── 知性は能力でも所有物でもなく、構造的条件下で成立する力学的状態である ── と、明らかに別系統の研究路線から収束している。この収束の意味は \S{}5 および \S{}10.2 で詳しく扱う。

\subsection{Contributions and scope}


本論文の主たる貢献は以下である:

\begin{enumerate}
\item \textbf{層分離の提示}(\S{}2):知覚・認識・知識・知性が異なる構造層に属することの構造論的分離
\item \textbf{三公理の定式化}(\S{}3):知性的動態が成立するための最小公理 ── 連結状態空間、時間順序付け、グローバル制約 $\Phi$ ── の同時必要性
\item \textbf{語彙基盤の確立}(\S{}4):IFGT を規範性なき記述語彙として位置づけ、Tero-Kobayashi-Nakagaki モデルとの語彙対応を提示
\item \textbf{領域横断的構造比較}(\S{}5):粘菌と Transformer の構造対応による、三公理の substrate 非依存性の検討
\item \textbf{Gettier 問題の再定位}(\S{}6):層間ミスマッチを検出する診断装置としての再解釈
\item \textbf{正しさの不安定化定理}(\S{}7):正しい知識の累積が分岐爆発を生む構造的帰結の動力学的記述
\item \textbf{個体内非閉論証}(\S{}8):因果推論の再帰構造から、個体内に停止演算子が存在しないことの構造的論証
\item \textbf{境界の命名と層対応}(\S{}9):\emph{the differential horizon of intelligence} の命名、および物理層・生命層・意識層・知性層の四層対応の提示
\end{enumerate}

本論文の \textbf{射程外} は明示的に以下を含む:

\begin{itemize}
\item \textbf{社会制度・集合的意思決定・ガバナンス}:本論文は越境点の構造的必要性のみを定式化し、外化された記述窓で展開される具体的構造(Part II 以降の課題)は扱わない
\item \textbf{倫理・規範・方法論的処方}:本論文は記述的構造論であり、規範的提言を一切含まない
\item \textbf{実装プロトコル・性能評価}:本論文は工学的指針ではなく構造的制約を提示する
\item \textbf{意識・主観性・経験的内容}:これらは別途の構造的扱いを要する領域であり、本論文の三公理は意識の有無を一切前提しない
\item \textbf{越境後に外化される具体構造}:命名と構造的必要性の論証は本論文で行うが、log、合意、制度、分散的改訂可能性などの具体的記述は別論文に譲る
\end{itemize}

本論文は \textbf{答えではなく制約を定式化する}。提示される六つの制約(\S{}10.1 に集約)は説明ではなく必要条件であり、これを満たさない知性論は、本論文が示す不安定性を再生産する可能性が高い。
